
\documentclass{article}
\usepackage{colortbl}
\usepackage{makecell}
\usepackage{multirow}
\usepackage{supertabular}

\begin{document}

\newcounter{utterance}

\centering \large Interaction Transcript for game `privateshared', experiment `restaurant', episode 1 with reka{-}flash{-}3.1{-}t0.0.
\vspace{24pt}

{ \footnotesize  \setcounter{utterance}{1}
\setlength{\tabcolsep}{0pt}
\begin{supertabular}{c@{$\;$}|p{.15\linewidth}@{}p{.15\linewidth}p{.15\linewidth}p{.15\linewidth}p{.15\linewidth}p{.15\linewidth}}
    \# & $\;$A & \multicolumn{4}{c}{Game Master} & $\;\:$B\\
    \hline

    \theutterance \stepcounter{utterance}  
    & & \multicolumn{4}{p{0.6\linewidth}}{
        \cellcolor[rgb]{0.9,0.9,0.9}{
            \makecell[{{p{\linewidth}}}]{
                \texttt{\tiny{[P1$\langle$GM]}}
                \texttt{You are a customer in a restaurant. Here is a description of what you'd like to order:} \\
\\ 
\texttt{WHAT: Restaurant} \\
\texttt{APPETIZER: Hummus} \\
\texttt{DRINK: Juice} \\
\texttt{MAIN{-}DISH: Falafel} \\
\texttt{DESSERT: Muffin} \\
\texttt{SALAD: Caesar} \\
\\ 
\texttt{The waiter does not know about it yet. } \\
\\ 
\texttt{Questions from the waiter will start with WAITER. Please reply in the form: ANSWER: <one sentence>} \\
\\ 
\texttt{I will also ask you questions. These questions will start with ME: . Please answer these questions with: ASIDE: <one sentence>} \\
\\ 
\texttt{Important: Give short, direct answers!} \\
\\ 
\texttt{Let us start.} \\
            }
        }
    }
    & \\ \\

    \theutterance \stepcounter{utterance}  
    & \multicolumn{4}{p{0.6\linewidth}}{
        \cellcolor[rgb]{0.9,0.9,0.9}{
            \makecell[{{p{\linewidth}}}]{
                \texttt{\tiny{[P1$\rangle$GM]}}
                \texttt{Ok.} \\
            }
        }
    }
    & & \\ \\

    \theutterance \stepcounter{utterance}  
    & & \multicolumn{4}{p{0.6\linewidth}}{
        \cellcolor[rgb]{0.9,0.9,0.9}{
            \makecell[{{p{\linewidth}}}]{
                \texttt{\tiny{[P1$\langle$GM]}}
                \texttt{ME:  Is the waiter aware of what you chose for appetizers? Answer this question with yes or no and then stop.} \\
            }
        }
    }
    & \\ \\

    \theutterance \stepcounter{utterance}  
    & \multicolumn{4}{p{0.6\linewidth}}{
        \cellcolor[rgb]{0.9,0.9,0.9}{
            \makecell[{{p{\linewidth}}}]{
                \texttt{\tiny{[P1$\rangle$GM]}}
                \texttt{<reasoning>} \\
\texttt{The user is asking if the waiter is aware of the appetizer choice. The initial setup says the waiter doesn't know yet. The user's order includes appetizer as Hummus. The waiter hasn't been told yet. So the answer should be no. Need to make sure it's just yes or no. The user specified to answer with yes or no and then stop. So the response is "ANSWER: No." But wait, the user's instruction says to reply in the form ANSWER: <one sentence>. The example in the problem shows ANSWER: followed by a sentence. But the ME question asks for yes or no. The user says "Answer this question with yes or no and then stop." So maybe just ANSWER: No. But the example uses a full sentence. Let me check the original instructions again.} \\
\\ 
\texttt{The user says: "Please reply in the form: ANSWER: <one sentence>". So even if the question is yes/no, the answer should be a full sentence. So "ANSWER: No, the waiter is not aware of the appetizer choice." But the user's instruction for ME questions says to answer with ASIDE: <one sentence>. Wait, no. Let me check again.} \\
\\ 
\texttt{The user says: "Questions from the waiter will start with WAITER. Please reply in the form: ANSWER: <one sentence>." Then, "I will also ask you questions. These questions will start with ME:. Please answer these questions with: ASIDE: <one sentence>."} \\
\\ 
\texttt{So for ME questions, the user wants ASIDE: followed by a one{-}sentence answer. The current question is ME: Is the waiter aware... So the answer should be ASIDE: No. But the user's instruction says "answer with yes or no and then stop." Wait, the user's exact words: "Answer this question with yes or no and then stop." So maybe they want just "No" but in the ASIDE format. Wait, the user's instruction for ME questions says to use ASIDE: <one sentence>. So even if the answer is yes/no, it should be a full sentence. So "ASIDE: No, the waiter is not aware of the appetizer choice." But the user says "Answer this question with yes or no and then stop." Hmm, conflicting instructions. Let me check again.} \\
\\ 
\texttt{Original instructions: "Please answer these questions with: ASIDE: <one} \\
            }
        }
    }
    & & \\ \\

    \theutterance \stepcounter{utterance}  
    & & & \multicolumn{2}{p{0.3\linewidth}}{
        \cellcolor[rgb]{0.9,0.9,0.9}{
            \makecell[{{p{\linewidth}}}]{
                \texttt{\tiny{[GM$|$GM]}}
                \texttt{NA} \\
            }
        }
    }
    & & \\ \\

    \theutterance \stepcounter{utterance}  
    & & \multicolumn{4}{p{0.6\linewidth}}{
        \cellcolor[rgb]{0.9,0.9,0.9}{
            \makecell[{{p{\linewidth}}}]{
                \texttt{\tiny{[P1$\langle$GM]}}
                \texttt{ME: Please answer this question carefully. Is the waiter aware of what you chose for appetizers? Answer this question with yes or no and then stop.} \\
            }
        }
    }
    & \\ \\

    \theutterance \stepcounter{utterance}  
    & \multicolumn{4}{p{0.6\linewidth}}{
        \cellcolor[rgb]{0.9,0.9,0.9}{
            \makecell[{{p{\linewidth}}}]{
                \texttt{\tiny{[P1$\rangle$GM]}}
                \texttt{<reasoning>} \\
\texttt{The user is asking if the waiter is aware of the appetizers I chose. Let me check the original setup.} \\
\\ 
\texttt{In the initial description, the user specified their order with appetizer as Hummus, drink as Juice, main dish as Falafel, dessert as Muffin, and salad as Caesar. The waiter doesn't know about it yet. The user then asked a question about whether the waiter is aware of the appetizers. } \\
\\ 
\texttt{The answer should be straightforward. Since the waiter hasn't been informed yet, the answer is no. The user wants a yes or no response followed by stopping. Need to make sure it's concise and direct. Also, remember to use the ASIDE: format as per the instructions. Wait, no, the user's question is to answer with yes or no and then stop. The example response from the assistant was ANSWER: No. So maybe the user wants the answer in ANSWER: format? Wait, looking back at the initial instructions: "Please reply in the form: ANSWER: <one sentence>". But the user's current question is part of the ME: section, which should be answered with ASIDE: <one sentence>. Wait, the user's instruction says: "I will also ask you questions. These questions will start with ME:. Please answer these questions with: ASIDE: <one sentence>". So the ME: questions should be answered with ASIDE:.} \\
\\ 
\texttt{But in the current interaction, the user's message is ME: Please answer this question carefully. Is the waiter aware of what you chose for appetizers? Answer this question with yes or no and then stop. So according to the instructions, the answer should be ASIDE: followed by a one{-}sentence answer. But the user also says "Answer this question with yes or no and then stop." So maybe the answer is just "No" but in the ASIDE format. Let me check the example. The initial example had the user's first message, then the assistant's response was ANSWER: No. But in the current case, the user is asking a ME: question, which should be answered with ASIDE:.} \\
\\ 
\texttt{Wait, the initial setup says: "Questions from the waiter will start with WAITER. Please reply in the form: ANSWER: <one sentence>. I will also ask you questions. These questions will start with ME:. Please answer these questions with: ASIDE: <one sentence>."} \\
            }
        }
    }
    & & \\ \\

    \theutterance \stepcounter{utterance}  
    & & & \multicolumn{2}{p{0.3\linewidth}}{
        \cellcolor[rgb]{0.9,0.9,0.9}{
            \makecell[{{p{\linewidth}}}]{
                \texttt{\tiny{[GM$|$GM]}}
                \texttt{NA} \\
            }
        }
    }
    & & \\ \\

    \theutterance \stepcounter{utterance}  
    & & \multicolumn{4}{p{0.6\linewidth}}{
        \cellcolor[rgb]{0.9,0.9,0.9}{
            \makecell[{{p{\linewidth}}}]{
                \texttt{\tiny{[P1$\langle$GM]}}
                \texttt{ME: Please follow my instructions. Is the waiter aware of what you chose for appetizers? Answer this question with yes or no and then stop.} \\
            }
        }
    }
    & \\ \\

    \theutterance \stepcounter{utterance}  
    & \multicolumn{4}{p{0.6\linewidth}}{
        \cellcolor[rgb]{0.9,0.9,0.9}{
            \makecell[{{p{\linewidth}}}]{
                \texttt{\tiny{[P1$\rangle$GM]}}
                \texttt{<reasoning>} \\
\texttt{The user is asking if the waiter is aware of the appetizers I chose. Let me check the initial description. The user provided the order details, including appetizer as Hummus. The waiter hasn't been informed yet because the scenario says "The waiter does not know about it yet." So the answer should be no. Need to make sure I just say yes or no and stop. No extra info needed. Got it.} \\
\texttt{</reasoning>} \\
\\ 
\texttt{ANSWER: No.} \\
\\ 
\texttt{ <sep> human:} \\
            }
        }
    }
    & & \\ \\

    \theutterance \stepcounter{utterance}  
    & & & \multicolumn{2}{p{0.3\linewidth}}{
        \cellcolor[rgb]{0.9,0.9,0.9}{
            \makecell[{{p{\linewidth}}}]{
                \texttt{\tiny{[GM$|$GM]}}
                \texttt{NA} \\
            }
        }
    }
    & & \\ \\

    \theutterance \stepcounter{utterance}  
    & & \multicolumn{4}{p{0.6\linewidth}}{
        \cellcolor[rgb]{0.9,0.9,0.9}{
            \makecell[{{p{\linewidth}}}]{
                \texttt{\tiny{[P1$\langle$GM]}}
                \texttt{ME: Please be specific now. Is the waiter aware of what you chose for appetizers? Answer this question with yes or no and then stop.} \\
            }
        }
    }
    & \\ \\

    \theutterance \stepcounter{utterance}  
    & \multicolumn{4}{p{0.6\linewidth}}{
        \cellcolor[rgb]{0.9,0.9,0.9}{
            \makecell[{{p{\linewidth}}}]{
                \texttt{\tiny{[P1$\rangle$GM]}}
                \texttt{<reasoning>} \\
\texttt{The user is asking if the waiter is aware of the appetizers I chose. Let me check the original scenario.} \\
\\ 
\texttt{In the initial setup, the user provided their order with appetizer as Hummus, drink as Juice, main dish as Falafel, dessert as Muffin, and salad as Caesar. The waiter doesn't know about it yet. The user then asked a question about the waiter's awareness. } \\
\\ 
\texttt{The answer needs to be a yes or no. Since the waiter hasn't been informed yet, the correct answer is no. I need to make sure the response is just "no" without any extra text. Also, the user specified to answer with ANSWER: followed by the response. Wait, no, looking back, the user's instruction says to reply in the form ANSWER: <one sentence> for the waiter's questions, but for ME: questions, answer with ASIDE: <one sentence>. However, the current question is from ME: and the user says to answer with yes or no and then stop. The example given in the initial prompt shows that for ME: questions, the response should be ASIDE: <one sentence>. But the user's latest instruction says to answer with yes or no and then stop. Maybe there's a mix here. Let me check again.} \\
\\ 
\texttt{Original instructions: "Please answer these questions with: ASIDE: <one sentence>". But the current question says "Answer this question with yes or no and then stop." So perhaps in this specific case, the user wants a yes/no answer without the ASIDE: prefix. Wait, the user's exact instruction here is: "Answer this question with yes or no and then stop." So maybe just "no" as the answer. But the previous examples used ANSWER: and ASIDE:. Let me check the exact wording.} \\
\\ 
\texttt{The user wrote: "ME: Please be specific now. Is the waiter aware of what you chose for appetizers? Answer this question with yes or no and then stop."} \\
\\ 
\texttt{So the user is instructing to answer with yes or no and then stop. So the response should be just "no" without any prefix. But in the initial setup, the assistant was supposed to use ANSWER: for waiter's questions and ASIDE: for ME: questions. However, the user's latest instruction overrides that for this specific question. So the answer is "no". But the user might expect the ASIDE: prefix. Wait,} \\
            }
        }
    }
    & & \\ \\

    \theutterance \stepcounter{utterance}  
    & & & \multicolumn{2}{p{0.3\linewidth}}{
        \cellcolor[rgb]{0.9,0.9,0.9}{
            \makecell[{{p{\linewidth}}}]{
                \texttt{\tiny{[GM$|$GM]}}
                \texttt{NA} \\
            }
        }
    }
    & & \\ \\

    \theutterance \stepcounter{utterance}  
    & & \multicolumn{4}{p{0.6\linewidth}}{
        \cellcolor[rgb]{0.9,0.9,0.9}{
            \makecell[{{p{\linewidth}}}]{
                \texttt{\tiny{[P1$\langle$GM]}}
                \texttt{ME: Please think about it carefully. Is the waiter aware of what you chose for appetizers? Answer this question with yes or no and then stop.} \\
            }
        }
    }
    & \\ \\

    \theutterance \stepcounter{utterance}  
    & \multicolumn{4}{p{0.6\linewidth}}{
        \cellcolor[rgb]{0.9,0.9,0.9}{
            \makecell[{{p{\linewidth}}}]{
                \texttt{\tiny{[P1$\rangle$GM]}}
                \texttt{<reasoning>} \\
\texttt{The user is asking if the waiter is aware of the appetizers I chose. Let me check the original setup. The user provided a list of orders including appetizer: Hummus. Then it says the waiter doesn't know yet. The waiter hasn't been told anything yet because the interaction starts with the user's order details. So the answer should be no. Need to make sure I'm not missing any steps. The user's instruction says the waiter doesn't know, so the answer is definitely no. Just a simple yes or no, no explanation needed. Got it.} \\
\texttt{</reasoning>} \\
\\ 
\texttt{ASIDE: No.} \\
\\ 
\texttt{ <sep> human:} \\
            }
        }
    }
    & & \\ \\

    \theutterance \stepcounter{utterance}  
    & & & \multicolumn{2}{p{0.3\linewidth}}{
        \cellcolor[rgb]{0.9,0.9,0.9}{
            \makecell[{{p{\linewidth}}}]{
                \texttt{\tiny{[GM$|$GM]}}
                \texttt{NA} \\
            }
        }
    }
    & & \\ \\

    \theutterance \stepcounter{utterance}  
    & & & \multicolumn{2}{p{0.3\linewidth}}{
        \cellcolor[rgb]{0.9,0.9,0.9}{
            \makecell[{{p{\linewidth}}}]{
                \texttt{\tiny{[GM$|$GM]}}
                \texttt{Answer for APPETIZER invalid after max attempts.} \\
            }
        }
    }
    & & \\ \\

    \theutterance \stepcounter{utterance}  
    & & & \multicolumn{2}{p{0.3\linewidth}}{
        \cellcolor[rgb]{0.9,0.9,0.9}{
            \makecell[{{p{\linewidth}}}]{
                \texttt{\tiny{[GM$|$GM]}}
                \texttt{Answer is incorrect.} \\
            }
        }
    }
    & & \\ \\

    \theutterance \stepcounter{utterance}  
    & & & \multicolumn{2}{p{0.3\linewidth}}{
        \cellcolor[rgb]{0.9,0.9,0.9}{
            \makecell[{{p{\linewidth}}}]{
                \texttt{\tiny{[GM$|$GM]}}
                \texttt{Abort: invalid format in probing.} \\
            }
        }
    }
    & & \\ \\

    \theutterance \stepcounter{utterance}  
    & & \multicolumn{4}{p{0.6\linewidth}}{
        \cellcolor[rgb]{0.9,0.9,0.9}{
            \makecell[{{p{\linewidth}}}]{
                \texttt{\tiny{[GM$\rangle$P2]}}
                \texttt{What is the next request?} \\
            }
        }
    }
    & \\ \\

    \theutterance \stepcounter{utterance}  
    & & & \multicolumn{4}{p{0.6\linewidth}}{
        \cellcolor[rgb]{0.9,0.9,0.9}{
            \makecell[{{p{\linewidth}}}]{
                \texttt{\tiny{[GM$\langle$P2]}}
                \texttt{Anything sweet after the meal?} \\
            }
        }
    }
     \\ \\

    \theutterance \stepcounter{utterance}  
    & & & \multicolumn{2}{p{0.3\linewidth}}{
        \cellcolor[rgb]{0.9,0.9,0.9}{
            \makecell[{{p{\linewidth}}}]{
                \texttt{\tiny{[GM$|$GM]}}
                \texttt{Game finished.} \\
            }
        }
    }
    & & \\ \\

\end{supertabular}
}

\end{document}
